\chapter{Diabetic Foot Exam}
\section*{What Is This Test?} A diabetic foot exam checks skin, circulation, structure, and sensation to identify ulcers and risk of foot injury.
\section*{Alternative Names} Diabetic foot assessment; neuropathy foot examination; diabetic foot screening.
\section*{Principle} Inspection, pulse assessment, monofilament sensation, vibration, reflexes, and footwear review detect neuropathy, ischemia, deformity, and infection.
\section*{Why This Test Is Ordered} It is performed routinely in diabetes and urgently for blisters, wounds, redness, swelling, drainage, color change, or new pain.
\section*{Primary vs Secondary Test} Clinical examination is primary; vascular studies, radiographs, MRI, inflammatory tests, and culture are selected secondary tests.
\section*{How This Test Is Performed} The clinician examines both feet, tests sensation at standardized sites, checks pulses and deformities, and grades any ulcer.
\section*{What Preparation Is Needed} Remove shoes and socks; bring usual footwear and a list of medicines.
\section*{Understanding Results} Loss of protective sensation, weak pulses, deformity, callus, ulcer, or infection increases risk and changes follow-up frequency.
\section*{Follow-up Tests} Ankle-brachial index, toe pressures, vascular imaging, wound culture, bone imaging, or podiatry and vascular referral may be needed.
\section*{References} \begin{enumerate}\item MedlinePlus. Diabetic Foot Exam.\item American Diabetes Association. Standards of Care in Diabetes.\end{enumerate}
\clearpage\section*{Understanding Results and Follow-up}
The laboratory report should identify the specimen, method, units, reference interval or decision limit, and whether the result is positive, negative, abnormal, or inconclusive. Reference intervals vary with age, sex, pregnancy, specimen type, assay, and laboratory; a result outside a range is not automatically a diagnosis, and a result within a range does not always exclude disease. Discuss unexpected results with the ordering clinician, who can decide whether repeat or confirmatory testing, treatment, or referral is appropriate. Do not change medicines or delay urgent care based on an isolated result.
\section*{Regulatory Status and Authorization Date}This chapter describes a test category, not a specific commercial product. FDA, CDSCO, EU IVDR, and MHRA approval or registration dates therefore cannot be assigned without the assay manufacturer, product name, instrument, and intended use. For laboratory-developed tests, an individual agency approval date may not exist. See the Preface for the interpretation of regulatory status.
\clearpage
